% ============================================================
\section{Experiment 4: Does the Screen Work Prospectively?}
\label{sec:exp4}
% ============================================================

\label{sec:analysis:heldout}
% The mechanism paragraph that used to follow this one -- the decomposition into the frozen and adapted
% arms and the exchangeability ratio -- now sits beside Table \ref{tab:heldout_decomp} in the appendix,
% which is the table it describes and which had no prose of its own. The finding itself stays here, and
% it is also one of the two lessons in the box at the end of S3.
\textbf{Held-out scoring is a precondition, not a refinement.}
Every number in this paper is scored on windows disjoint from the ones used
for early stopping.  Scoring
$\denc$ on the selection windows instead \textbf{reverses its
sign in 4 of the 31 intervention cells}, and \emph{every one of the four
reverses in the same direction}: from ``the frozen encoder was better'' to
``full fine-tuning was better'' (Appendix Table~\ref{tab:heldout_decomp}).  The largest
is Chronos/ETTh1, which moves from $+6.8$ to $-39.2$, a 46-point swing.  Two
of the four are among the seven gate-passing cells, so this is not confined to
cells nobody would have looked at, and the reversals are driven by the frozen arm,
which a distribution shift between the splits penalises more heavily than the
adapted arm because it cannot adapt to one.

\label{sec:analysis:prospective}
\textbf{A pre-registered prospective test with no positive instance.}
Everything above is retrospective: the cells were run, then read, and a
diagnostic that only explains cells it has seen is worth little, so we tested
ours forward.  \textbf{Predictions were written to
\texttt{results/v47\_prospective/preregistration.json} and committed to git at
a named revision before any outcome run for those cells existed}; the scoring
script reads that file and the finished runs and decides nothing.
% The two registered rules were an itemize; folded into one sentence for page room, with nothing
% dropped -- both rules, the pre-correction denominator and the reason for registering a competitor
% are all still stated. The itemize cost four printed lines of its own.
Two rules were registered against each other: the \textbf{gate rule},
gate$_\text{pre}\!\geq\!0.20 \Rightarrow$ at risk, which is the paper's own
criterion frozen \emph{before} the baseline correction and therefore graded on
trend-denominator scores rather than re-derived after seeing the outcomes; and a
\textbf{dataset rule}, dataset ${\in}$ \{ETTh1, Weather\} $\Rightarrow$
degradation, a \emph{post hoc} pattern read off the original Moirai cells and
registered as a competitor so the gate is graded against something.
Eight cells were then run (Weather, ETTm2, Electricity7, ETTh1
$\times$ Small/Base/Large, 3 seeds each) and scored by the pre-registered
outcome definition, unchanged (Appendix Table~\ref{tab:prospective}).
% The eight-row per-cell table moved to the appendix, beside the threshold sweep over the same eight
% cells, which is the table it is read against. It was 14 printed lines of a 9-page body and every
% claim the body makes from it -- TP 0, FP 8, the one cell that clears a corrected gate, the seed
% counts -- is stated in the two paragraphs below and registered in check_paper_numbers.py.

\textbf{Result: the outcome the gate was registered to predict never occurs.}
None of the eight cells degraded.  The gate flagged all eight, giving
\textbf{TP~0, FP~8, FN~0}: precision $0.00$, while \textbf{sensitivity and
specificity are both unestimable}---sensitivity has an empty positive class in its
denominator, and specificity is $0/8$ only because every prediction was positive.
The competing dataset rule scores no better, for the same reason.  \textbf{What
``prospective'' does and does not mean here:} the predictions were frozen before
these cells were \emph{trained}, so they are prospective with respect to the
intervention---but the datasets and their splits already existed, so this is not a
test on newly collected data.  Nor is the empty positive class an artefact of the
corrected gate: only one of the eight clears $0.20$ on either corrected split
(base/ETTm2 $h{=}192$, $+0.21$ on selection windows) and it fails clause~(ii) in 0
of 3 seeds (Appendix~\ref{app:corrections}).  \textbf{Nor is the $0.20$ threshold
what fails}: re-scoring the same predictor against the same outcome at thresholds
from $0.10$ to $0.60$, moving nothing else, gives precision $0.00$ at every one,
with 8, 8, 6, 5, 3 and 2 cells flagged (Appendix
Table~\ref{tab:gate_sensitivity})---the problem is not the cutoff but that the
predicted class is empty.

\textbf{The gate does not provide a reliable ordering, and there is no actionable
threshold.}
Within Moirai ($n{=}22$ cells from six series, the only arm with enough cells to
ask), the gate's continuous score correlates \emph{negatively} with held-out
$\denc$: Spearman $\rho{=}{-}0.293$, and neither cells nor series exclude zero, so
we report the direction and not the interval (Appendix~\ref{app:clustered}).
The sign is the wrong one---higher measured pre-trained value goes with encoder
adaptation helping \emph{more}, not less, the opposite of what a preservation
heuristic needs---but the sharper point is that the magnitude is not stable:
scoring the same gate on the held-out windows instead gives $\rho{=}{-}0.518$
on the same cells.
An ordering that nearly doubles when the gate changes split is not one a
practitioner can set a threshold on, and no other quantity we measured orders
$\denc$ either (Appendix~\ref{app:predictor}).

\textbf{What survives.} Not the gate as a predictor of harm, nor as a screen for
when preservation is worth the effort---here it fired eight times and was wrong
eight times---but the gate as a \emph{disqualifier}: a cell that fails it cannot
supply evidence about damage to pre-trained features, and 24 of our 31 fail it.
